% Source-aligned English for printed p.181.
\SourcePage{181}

\LayerTarget{tohoku:proof:proposition-3.9.2:finite-intersection:p2}
The $O_x(\mathbf R^{L_x})$ form an open covering of $X$, so there is a finite
subset $Y\subset X$ such that the $O_x(\mathbf R^{L_x})$ corresponding to
$x\in Y$ cover $X$. Let $L$ be the union of the $L_x$ for $x\in Y$. I claim
that the covering $\mathbf R^L$ of $X$ refines $\mathcal U$. Indeed, let
$A\in\mathbf R^L$ be nonempty and let $a\in A$. There is an $x\in Y$ such
that $a\in O_x(\mathbf R^{L_x})$, whence
$E_a(\mathbf R^{L_x})\subset E_x(\mathbf R^{L_x})\subset U_x$ and, a
fortiori,
$A\subset E_a(\mathbf R^L)\subset E_a(\mathbf R^{L_x})$ is contained in
$U_x$. This proves the assertion. Applying this result, for every
$A\in\mathcal S$, to the family of coverings of $A$ induced by the coverings
$\mathbf R\in\mathfrak R(X)$ gives the desired conclusion.

\LayerTarget{tohoku:application:corollary-3.9.1-3:cubes}
The most striking application of Corollary 3 is the case in which $X$ is the
compact cube $0\leq x_i\leq1$ in $\mathbf R^n$, $\mathcal S$ is the family
of compact cubes $A$ of the form $a_i\leq x_i\leq b_i$ contained in $X$,
and $\mathfrak R(A)$ is the family of coverings of $A$ by two sets defined by
hyperplanes parallel to the coordinate hyperplanes. To verify that
$H^i(A,F)=0$ for $i>0$ and every $A$, it is enough to verify that, for every
$A$ and every section $f$ of $F$ on $A_1\cap A_2$, one has $f=f_1-f_2$,
where $f_i$ is a section of $F$ on $A_i$. This is the reduction made by
H. Cartan in his proof of the fundamental theorems on Stein manifolds [5].

\LayerTarget{tohoku:remark:theorem-3.9.1:nonabelian-n-one}
\noindent\textsc{Remark.} If $n=1$, Theorem 3.9.1 remains meaningful and is
easily verified directly when $F$ is assumed to be a sheaf of groups that
need not be abelian. This simplifies the proof of theorem [5, XVIII] on
invertible holomorphic matrices.

\LayerTarget{tohoku:section:3.10}
\SectionHeading{3.10. Passage to the limit in sheaf cohomology.}
We shall give only two results in this direction, one of which will be used
in Chapter 5, No. 7. They are special cases of the following general result
of homological algebra.

\LayerTarget{tohoku:proposition:3.10.1}
\noindent\textsc{Proposition 3.10.1.} \textit{Let $\mathcal C,\mathcal C'$ be
two abelian categories. Suppose that every object of $\mathcal C$ is
isomorphic to a subobject of an injective object, and that $\mathcal C'$
satisfies axiom AB 5 (cf. 1.5), which in particular permits the formation of
inductive limits in $\mathcal C'$ (cf. Proposition 1.8). Let
$(F_i)_{i\in I}$ be an inductive system of covariant additive functors from
$\mathcal C$ to $\mathcal C'$, and let $F=\varinjlim F_i$ be their inductive
limit, defined by $F(A)=\varinjlim F_i(A)$ for every $A\in\mathcal C$. The
homomorphisms $F_i\longrightarrow F$ induce homomorphisms of
$\partial$-functors $(\mathbf R^pF_i)\longrightarrow(\mathbf R^pF)$, and
hence a homomorphism of $\partial$-functors:}

\LayerTarget{tohoku:eq:3.10.1}
\[
\tag{3.10.1}
\varinjlim_i \mathbf R^pF_i(A)\longrightarrow \mathbf R^pF(A).
\]
\textit{(The connecting homomorphisms for the sequence of functors
$\varinjlim_i\mathbf R^pF_i$ are defined as the inductive limits of the
connecting homomorphisms for the $\mathbf R^pF_i$.) The homomorphisms
(3.10.1) are isomorphisms.}

\LayerTarget{tohoku:proof:proposition-3.10.1:p1}
It is enough to take an injective resolution $C=C(A)$ of $A$. The first
member of (3.10.1) is then $\varinjlim_i H^p(F_iC(A))$, while the second is
$H^p(\varinjlim_i F_iC(A))$. They are isomorphic because the functor
$\varinjlim_i$ on the category of inductive systems over $I$ with values in
$\mathcal C'$ is exact (Proposition 1.8) and, in particular, commutes with
taking the homology of complexes.

\TohokuContinuation{182}
